% Level 3: laws, and the Lean they need
\begin{frame}
  \frametitle{Selections Move Down, Under Side Conditions}
  {\centering\renewcommand\arraystretch{1.25}
    {\headerrow\begin{tabular}{@{}r@{\quad}l@{\qquad}l@{}}
      \hline
      & \textbf{Law} & \textbf{Holds when} \\ \hline
      1 & $\sel{p}(\sel{q}(R)) = \sel{p \wedge q}(R)$ & always \\
      2 & $\sel{p}(\sel{q}(R)) = \sel{q}(\sel{p}(R))$ & always \\
      3 & $\sel{p}(R \cup S) = \sel{p}(R) \cup \sel{p}(S)$ & always \\
      4 & $\sel{p}(R - S) = \sel{p}(R) - S$ & always \\
      5 & $\sel{p}(R \times S) = \sel{p}(R) \times S$ & $p$ mentions only
        attributes of $R$ \\
      6 & $\sel{p}(R \join S) = \sel{p}(R) \join S$ & $p$ mentions only
        attributes of $R$ \\ \hline
    \end{tabular}}\par}
  \medskip
  In Lean, each side condition is a \textbf{hypothesis} of the theorem.
  Without it, law 5 is false: $\sel{\att{week} = 2}$ cannot move onto
  $\rel{Member}$, which has no $\att{week}$.
  \slideref{\lnk{ahv6} (section 6.1)}
  \leveltag{3}
\end{frame}

\begin{frame}
  \frametitle{Projections and Joins Have Laws, and Traps}
  {\centering\renewcommand\arraystretch{1.25}
    {\headerrow\begin{tabular}{@{}r@{\quad}l@{\qquad}l@{}}
      \hline
      & \textbf{Law} & \textbf{Holds when} \\ \hline
      7 & $\proj{A}(\proj{B}(R)) = \proj{A}(R)$ & $A \subseteq B$ \\
      8 & $\proj{A}(R \cup S) = \proj{A}(R) \cup \proj{A}(S)$ & always \\
      9 & $\proj{A}(\sel{p}(R)) = \sel{p}(\proj{A}(R))$ & $p$ mentions
        only attributes in $A$ \\
      10 & $R \join S = S \join R$ & attributes are named \\
      11 & $(R \join S) \join T = R \join (S \join T)$ & always \\ \hline
    \end{tabular}}\par}
  \medskip
  \textbf{Not} a law: $\proj{A}(R - S) = \proj{A}(R) - \proj{A}(S)$.
  \textbf{Try}: find a counterexample.
  \slideref{\lnk{ahv6} (section 6.1)}
  \leveltag{3}
\end{frame}

\begin{frame}
  \frametitle{Lists, Bags or Sets: Laws Need the Right Equality}
  {\centering\renewcommand\arraystretch{1.3}
    {\headerrow\begin{tabular}{@{}l@{\qquad}c@{\qquad}c@{\qquad}l@{}}
      \hline
      & \textbf{Order matters} & \textbf{Duplicates matter} & \textbf{In
        Lean} \\ \hline
      list & yes & yes & \code{List} \\
      bag & no & yes & \code{Multiset} \\
      set & no & no & \code{Finset} \\ \hline
    \end{tabular}}\par}
  \begin{itemize}
    \item Law 1 holds even for lists; law 10 only up to reordering;
      $R \cup R = R$ only for sets.
    \item Relational algebra is about sets: start with lists if you like,
      but the laws need set equality.
  \end{itemize}
  \slideref{\lnk{mil-finset} \rfsep \lnk{ml-finset} \rfsep \lnk{ml-multiset}}
  \leveltag{3}
\end{frame}

\leanframe{Level 3 Needs Propositions, Tactics and Sets}{3}{%
  read a proposition as a type, and a proof as a program &
    \rds{fpil-props, tpil-pat} \\
  prove equalities with \code{rfl}, \code{rw} and \code{simp} &
    \rds{tpil-rw, tpil-simp} \\
  use hypotheses: \code{intro}, \code{exact}, \code{apply}, \code{cases} &
    \rds{tpil-basic} \\
  prove facts about lists by \code{induction} & \rds{fpil-tactics} \\
  prove two sets equal, element by element & \rds{mil-sets, mil-finset} \\
}{Practise induction and rewriting with the Natural Number Game, on the
  \lnk{games}.}
